\section{Latency and Memory Methodology}
\label{app:latency-methodology}

All timings and memory figures in this paper are measured, not derived from parameter or
FLOP counts. This appendix records the harness in full so the numbers can be reproduced or
contested. The harness is \texttt{experiments/run\_latency.py}; it times the native model
and every truncated variant in one process with an identical protocol.

\paragraph{The truncated model is real.}
Every layer-wise result elsewhere in this paper is obtained by running the full 12-block
encoder and reading intermediate states off forward hooks. That is exact for accuracy,
because the encoder is a feed-forward stack and blocks after $\ell$ cannot influence block
$\ell$'s output, but it means no truncated model is ever instantiated. For timing we build
one: encoder blocks $\ell+1 \ldots 12$ are deleted from the module list and the linear
adapter is spliced in front of the frozen final RMSNorm, which the encoder applies after
its last remaining block. Before timing, the harness asserts that the truncated encoder
reproduces the hooked full-model states; the observed maximum absolute deviation was
exactly 0.0e+00 across all measured depths, i.e.\ bit-identical.

\paragraph{Hardware.}
GPU: NVIDIA A100-SXM4-40GB with 39\,GiB of VRAM (compute capability 8.0), driver 580.178.04; 1 device visible, one used. Host CPU: AMD EPYC 7413 24-Core Processor, 2 cores allocated to the job.\footnote{Host CPU, core allocation and cuDNN version were recorded by a separate probe run (\texttt{latency\_\_env\_probe.json}) on the same GPU and job configuration; the timing run predates the harness capturing them. Every timing and memory figure in this appendix comes from the timing run itself.}

\paragraph{Software and precision.}
Python 3.11.4, PyTorch 2.12.1 (CUDA 13.2, cuDNN 92101 as reported by \texttt{torch.backends.cudnn.version()}), model \texttt{amazon/chronos-2}. 
All computation is in \texttt{float32} --- the precision at which every accuracy number
in this paper was produced. No \texttt{torch.compile}, no automatic mixed precision, and no
manual kernel fusion are applied (TF32 matmul: false; TF32 cuDNN: true).

\paragraph{Workload.}
Context $C=512$, horizon $H=64$, giving $K=4$ native forecast slots and an encoder sequence of 37 tokens (32 context patches, one REG token, and the forecast slots). Batch sizes 1, 32, 256 are reported: the smallest is the single-request
latency regime, the largest the batched-throughput regime, and they can support different
conclusions. Contexts are synthetic random walks; the forward pass has no data-dependent
control flow, so timing is independent of the data and needs no dataset.

\paragraph{Timing protocol.}
Each configuration is warmed up for 20 iterations, which are discarded, then
timed over 100 repetitions. CUDA is asynchronous, so
\texttt{torch.cuda.synchronize()} brackets every timed region; without it the measurement
records only Python dispatch. We report the \emph{median} over repetitions, which is robust
to occasional descheduling by the operating system, and the 95th percentile separately as
the tail. Two timings are recorded because they answer different questions:
\texttt{predict} calls \texttt{Chronos2Pipeline.predict\_quantiles} on host arrays and
therefore \emph{includes} preprocessing and the host-to-device transfer, i.e.\ the whole
serving path; \texttt{encode} times the encoder forward alone on a tensor already resident
on the device and therefore \emph{excludes} transfer. Speedups quoted in the main text use
\texttt{predict}, the conservative choice, since its fixed costs do not shrink with depth.

\paragraph{Throughput.}
Series per second is the batch size divided by the median \texttt{predict} time. It is not
the reciprocal of single-request latency: batching amortises fixed per-call cost, so
throughput improves with batch size while latency does not.

\paragraph{Memory.}
Peak memory is \texttt{torch.cuda.max\_memory\_allocated()} after
\texttt{reset\_peak\_memory\_stats()}, recorded over the timed region. We deliberately do
not use \texttt{nvidia-smi}, which reports the caching allocator's reservation and would
not shrink with truncation. The weight-only figure is the allocated total after loading and
truncating but before any forward pass. The model is reloaded from scratch for every depth
so that dropped blocks are genuinely freed rather than merely unused; a harness that reuses
one process would report a memory column that measures nothing.

\paragraph{Measured results.}
Table~\ref{tab:latency-full} reports every configuration. Latency and memory depend on
truncation depth only, not on the dataset, so the per-dataset speedups in the main text are
lookups into this table at each dataset's selected depth.

\begin{table}[t]
    \centering
    \small
    \caption{Measured latency, throughput and memory of the truncated model. \texttt{predict} is the full serving path (transfer included); \texttt{encode} is the encoder forward on a device-resident tensor (transfer excluded). Speedup is the median \texttt{predict} time of the full 12-block model divided by that of the truncated model at the same batch size.}
    \label{tab:latency-full}
    \begin{tabular}{llrrrrrr}
        \toprule
        Batch & Depth & Blocks & predict (ms) & p95 (ms) & encode (ms) & series/s & peak (MiB) \\
        \midrule
        1 & \(L1\) & 1 & 3.69 & 3.71 & 2.80 & 271 & 72.3 \\
         & \(L3\) & 3 & 6.49 & 6.58 & 5.61 & 154 & 144.5 \\
         & \(L6\) & 6 & 10.65 & 10.75 & 9.74 & 94 & 252.5 \\
         & \(L8\) & 8 & 13.43 & 13.58 & 12.52 & 74 & 324.5 \\
         & \(L10\) & 10 & 16.20 & 16.28 & 15.34 & 62 & 396.5 \\
         & \(L11\) & 11 & 17.60 & 17.73 & 16.72 & 57 & 432.5 \\
         & \(L12\) & 12 & 18.97 & 19.12 & 18.09 & 53 & 468.5 \\
        \midrule
        32 & \(L1\) & 1 & 5.42 & 5.98 & 3.52 & 5904 & 113.7 \\
         & \(L3\) & 3 & 9.23 & 9.31 & 7.33 & 3466 & 190.0 \\
         & \(L6\) & 6 & 14.91 & 15.21 & 13.02 & 2147 & 298.0 \\
         & \(L8\) & 8 & 18.67 & 18.88 & 16.79 & 1714 & 370.0 \\
         & \(L10\) & 10 & 22.46 & 22.58 & 20.57 & 1425 & 442.1 \\
         & \(L11\) & 11 & 24.36 & 24.39 & 22.48 & 1313 & 478.1 \\
         & \(L12\) & 12 & 26.29 & 26.34 & 24.39 & 1217 & 513.3 \\
        \midrule
        256 & \(L1\) & 1 & 28.09 & 28.66 & 18.30 & 9115 & 534.7 \\
         & \(L3\) & 3 & 55.45 & 56.06 & 45.70 & 4616 & 634.5 \\
         & \(L6\) & 6 & 96.58 & 97.13 & 86.86 & 2651 & 742.5 \\
         & \(L8\) & 8 & 123.97 & 124.49 & 114.42 & 2065 & 814.6 \\
         & \(L10\) & 10 & 151.91 & 152.63 & 142.64 & 1685 & 886.6 \\
         & \(L11\) & 11 & 165.20 & 165.86 & 155.66 & 1550 & 922.6 \\
         & \(L12\) & 12 & 178.95 & 179.45 & 169.02 & 1431 & 958.6 \\
        \bottomrule
    \end{tabular}
\end{table}
